
\section{Jacobian Spectral Radius Regularization}
Building on our previous experimental findings, we propose that effective test-time scalable latent reasoning must satisfy both effectiveness and stability. This requirement stems from the nature of reasoning as an iterative process of uncertainty reduction and thought refinement. This implies that hidden states should converge toward a stable and effective fixed point. Systems that are effective but unstable produce chaotic reasoning trajectories, whereas systems that are stable but ineffective yield shallow thoughts. To address this, we introduce STARS (STAbility-driven Recurrent Scaling), a training framework that combines Jacobian Spectral Radius Regularization (JSRR) with random loop sampling.
% 根据上一节的实验，我们认为/假设  test-time scalable la-
% tent reasoning must satisfy both effectiveness and stability. 这其中的道理在于reasoning is an iterative process of reducing uncertainty and refining thoughts. In terms of dynamics, this means the hidden states should converge toward a effective and stable fix point. If a system is effective but unstable, the thoughts become chaotic; if it is stable but ineffective, the thoughts remain shallow. we propose STARS (STAbility-driven Recurrent Scaling), a unified training framework that integrates Jacobian Spectral Radius Regularization (JSRR) with random loop sampling.

%citation
According to the Lyapunov Linearization Theorem for discrete-time dynamical systems, the local stability of a nonlinear system defined by $\mathbf{h}^{(t+1)} = \Phi_\theta(\mathbf{h}^{(t)})$ at a fixed point $\mathbf{h}^\star$ is governed by the properties of its Jacobian matrix, denoted as:$$J(\mathbf{h}^\star) = \left. \nabla_{\mathbf{h}} \Phi_\theta(\mathbf{h}) \right|_{\mathbf{h}=\mathbf{h}^\star}$$To determine stability, we examine the set of eigenvalues $\{\lambda_1, \lambda_2, \dots, \lambda_n\}$ of $J(\mathbf{h}^\star)$. The critical metric is the spectral radius, denoted by $\rho(J(\mathbf{h}^\star))$, which is defined as the maximum absolute value (modulus) of these eigenvalues:$$\rho(J(\mathbf{h}^\star)) = \max_{i} \{ |\lambda_i| \}$$Specifically, if the spectral radius satisfies: $\rho(J(\mathbf{h}^\star)) < 1$, the fixed point is \textit{asymptotically stable}. Under this condition, any small perturbations introduced to the system will decay exponentially over successive iterations. 
And if the spectral radius is smaller, the convergence rate also becomes faster.

\begin{table*}[t]
    \centering
    \caption{We report Accuracy (\%) on various mathematical benchmarks. Best results for LoopLMs are \textbf{bolded}.}
    \label{tab:ouro_iterations}
    
    \resizebox{0.99\textwidth}{!}{ 
    \begin{tabular}{l c c c c c c c}
        \toprule
        \textbf{Model} & \textbf{Recurrents} & \textbf{GSM8K} & \textbf{MATH500} & \textbf{ASDiv} & \textbf{SVAMP} & \textbf{AMC23}& \textbf{AVG}\\
        \midrule
        \multicolumn{8}{c}{\textit{Standard Language Models
        }} \\[0.5mm]
        Gemma3-1b-it~\citep{team2025gemma}  & - & 42.76 & 41.20 & 72.45 & 62.67 & 25.00 & 48.82\\
        \midrule
        Llama-3.2-3B-Instruct~\citep{grattafiori2024llama}
            & - & 71.11 &  42.40 &82.04 &  79.67 &  15.00 & 58.04\\
        \midrule
        Qwen2.5-3B-Instruct~\citep{team2024qwen2} 
             & - & 74.91& 68.20 & 76.18 & 87.00 & 32.50&  67.76\\
        \midrule
        Qwen3-1.7B~\citep{yang2025qwen3}
            & - & 75.82 & 65.80 & 85.64 & 89.00 &  17.50& 66.75\\
        \midrule
        Qwen3-4B~\citep{yang2025qwen3}
            & - & 80.36 & 63.00 & 86.59 & 89.00 &  12.50 & 66.29\\
        \toprule
        \multicolumn{8}{c}{\textit{Looped Language models
        }} \\[0.5mm]
        % \multirow{3}{*}{Recurrent-Llama-3.2-train-recurrence-8~\citep{mcleish2025teaching}}
        \multirow{3}{*}{Recurrent-Llama-3.2~\citep{mcleish2025teaching}}
            & 2 & 15.01 & 14.60 & 30.50 & 18.67 & 00.00 & 15.76\\ 
            & 4 & 14.03 & 14.60 & 41.65 & 18.67 & 5.000 & 18.79\\ 
            & 8 & 13.04 & 14.00 & 21.33 & 42.17 & 7.500 & 19.61\\
         \midrule
        % \multirow{3}{*}{Recurrent-OLMo-2-0425-train-recurrence-8~\citep{mcleish2025teaching}}  
        \multirow{3}{*}{Recurrent-OLMo-2-0425~\citep{mcleish2025teaching}} 
            & 2 & 17.89 & 9.000 & 29.33 & 31.97 & 00.00 & 17.64\\ 
            & 4 & 17.13 & 9.800 & 38.13 & 38.67 & 2.500 & 21.25\\ 
            & 8 & 17.97 & 11.40 & 40.00 & 38.48 & 00.00 & 21.57\\
         \midrule
        % ---------------- Group 1: Ouro ----------------
        \multirow{3}{*}{Ouro-1.4B~\citep{zhu2025scaling}} 
            & 2  & 76.88 & 54.80 & 74.27 & 81.67 & 35.00& 64.52\\
            &  4  & 75.21 & 59.60 & 76.57 & 75.67 & 50.00& 67.41\\
            % &  6  & 71.49 & 55.40 & 72.84 & 75.00 & 42.50\\
            &  8  & 58.23 & 40.80 & 70.07 & 66.33 & 40.00& 55.09\\
        \midrule
        
        % ---------------- Group 2: Ouro + SFT ----------------
        \multirow{3}{*}{Ouro-1.4B-SFT} 
            &  2  & 76.04 & 55.80 & 82.08 & 82.33 & 30.00& 65.25\\
            &  4  & 80.06 & 64.60 & 83.47 & 76.67 & 47.50&  70.46\\
            % &  6  & 67.93 & 45.80 & 79.61 & 77.00 & 32.50\\
            &  8  & 60.05 & 39.20 & 75.10 & 68.00 & 22.50& 52.97\\
        \midrule
        
        % ---------------- Group 3: Ouro + Method (Yours) ----------------

\multirow{3}{*}{Ouro-1.4B-STARS (Ours)}  
            &  2  & 75.89 & 56.40 & 82.56 & 79.67 & 40.00& 66.90\\
     &  4  & \textbf{81.96} & \textbf{67.40} & \textbf{84.73} & \textbf{84.33} & \textbf{52.50}& \textbf{74.18}\\
            % &  6  & 81.27 & 60.60 & 83.96 & 85.33 & 50.00\\
            &  8  & 74.45 & 54.80 & 82.52 & 81.00 & 35.00& 65.55\\
        % ---------------- Group 4: Qwen3-1.7B ----------------
        \bottomrule
    \end{tabular}}
\end{table*}

%cite
Previous works~\citep{bai2019deep,bai2021stabilizing}  often regularize using the Frobenius norm \(\|J\|_F\). However, while the spectral radius and the norm satisfy \(\rho(J) \le \|J\|\), directly constraining the norm is overly restrictive and may excessively compress the model's expressive capacity. In contrast, regularizing the spectral radius provides a mathematically precise means to achieve stability.

Direct eigenvalue computation of \(J \in \mathbb{R}^{D \times D}\) is infeasible for large \(D\). We therefore adopt an efficient spectral radius estimator based on the power iteration method. Starting from a randomly initialized vector \(\mathbf{v}\), the power iteration procedure repeatedly updates \(\mathbf{v} \leftarrow \frac{J\mathbf{v}}{\|J\mathbf{v}\|}\), eventually converging to the dominant eigenvector of \(J\). The corresponding spectral radius can then be estimated by: \(\rho(J) \approx \|J\mathbf{v}\|_2\).
To integrate this approach into large-scale model training, we use only  a single-step power iteration with the Jacobian-vector product technique in Pytorch, enabling efficient and memory-aware computation without explicit construction of the full Jacobian matrix.



We adopt a single-step update for two reasons: 1) multi-step iteration introduces complex gradient dependencies that can lead to abnormal gradients; 2) It is computationally lightweight, and experiments show it provides effective supervisory signals. Although the single-step estimate may be noisy for individual samples, its optimization direction remains statistically accurate across batches.

While the ultimate goal is to optimize the spectral radius at the fixed point $\mathbf{h}^\star$, identifying the exact location of the fixed point during training is computationally prohibitive. Consequently, we adopt a proxy approach: for a specific iteration $t$ within the LoopLM execution, we regulate the squared spectral radius of the Jacobian at the current state $\mathbf{h}^{(t)}$. For a batch of $N$ samples, the JSRR loss at iteration $t$ is formulated as:
% \begin{equation}
% \label{e2}
%     \mathcal{L}_{\text{JSRR}}^{(t)} = \frac{1}{N} \sum_{i=1}^N \left(\frac{\mathbf{v}_{1}^{(t,i)\top} J^{(t,i)} \mathbf{v}^{(t,i)}_{1}}{\mathbf{v}_{1}^{(t,i)\top} \mathbf{v}^{(t,i)}_{1}}\right)^2,
% \end{equation}
% \begin{equation}
% \label{e2}
%     \mathcal{L}_{\text{JSRR}}^{(t)} = \frac{1}{N} \sum_{i=1}^N \left(\frac{\mathbf{v}_{1}^{(t,i)\top} J^{(t,i)} \mathbf{v}^{(t,i)}_{1}}{\mathbf{v}_{1}^{(t,i)\top} \mathbf{v}^{(t,i)}_{1}}\right)^2,
% \end{equation}
\begin{equation}
\label{e2}
    \mathcal{L}_{\text{JSRR}}^{(t)} 
    = \frac{1}{N} \sum_{i=1}^N 
    \left\|
    J^{(t,i)} \mathbf{v}^{(t,i)}
    \right\|_2^2,
\end{equation}
where $J^{(t,i)} = \frac{\partial \mathbf{h}^{(t+1)}_{i}}{\partial \mathbf{h}^{(t)}_{i}}$ is the Jacobian of the $i$-th sample, and $\mathbf{v}_{1}^{(t,i)}$ is the one-step dominant eigenvector estimate of the $i$-th sample at iteration $t$.

However, directly constraining the spectral radius at only a single iteration \( t \) is suboptimal, as it neglects the stability properties across the entire inference trajectory. Therefore, we propose to integrate JSRR with the random loop sampling strategy introduced earlier. This ensures that spectral radius regularization is applied across a diverse set of states encountered during iterative inference, promoting global stability over the support set of the latent dynamics. Our final training objective is defined as the expectation of the combined loss over the recurrent depth distribution \( \mathcal{P} \):
\begin{equation}
    \mathcal{L}_{\text{STARS}} = \mathbb{E}_{t \sim \mathcal{P}} \left[ (1-\lambda) \cdot \mathcal{L}_{\text{SFT}}^{(t)} + \lambda \cdot \mathcal{L}_{\text{JSRR}}^{(t)} \right],
\end{equation}
where \(\lambda\) is a balancing hyperparameter. By sampling the loop length \(t\) from the distribution \(\mathcal{P}\), this formulation encourages the model to achieve high performance on the training data while maintaining a small spectral radius across the entire trajectory. 
The complete procedure of our algorithm is presented in Algorithm \ref{alg:stars_jacobian}.



